\chapter{Colonoscopy}

\section*{Overview}
A colonoscopy uses a long, flexible scope with a camera to examine the entire lining of the large intestine (colon) and rectum to detect inflammation, polyps, or cancer.

\section*{Alternative Names}
Lower GI endoscopy, Colonic evaluation

\section*{Principle}
A flexible endoscope with a camera is advanced through the anus into the colon, which is distended with gas to allow inspection of the entire large intestine. Instruments passed through the scope permit biopsy, polypectomy, and treatment of bleeding.
\section*{History}
Developed by William Wolff and Hiromi Shinya in 1969, introducing the technique of colonoscopic polypectomy which revolutionized colon cancer prevention.

\section*{Why This Test or Procedure Is Ordered}
Ordered for colorectal cancer screening starting at age 45, or to investigate chronic diarrhea, rectal bleeding, unexplained weight loss, or abdominal pain.

\section*{Is This a Primary or Secondary Test or Procedure}
Primary screening and preventive tool for colorectal cancer.

\section*{How This Test or Procedure Is Performed}
You lie on your side under sedation. The colonoscope is inserted through the rectum and slowly advanced to the cecum. Air or CO2 is used to inflate the colon.

\section*{What Preparation Is Needed}
Follow a clear liquid diet the day before and consume a strong laxative solution (bowel prep) to completely clear the colon of stool.

\section*{Contraindications and Precautions}
Toxic megacolon, fulminant colitis, or suspected bowel perforation.

\section*{Risks}
Bowel perforation (<1 in 1,000), bleeding post-polypectomy (<1 in 200), or sedation-related complications.

\section*{What Happens After the Test or Procedure}
You will remain in recovery until the sedative wears off. Mild gas and bloating are normal. Avoid driving for 24 hours.

\section*{Understanding Your Results}
If polyps are found, they are removed (polypectomy) and sent for biopsy. Normal results mean no polyps or lesions were detected.

\section*{Normal Results}
Clean mucosal lining of the colon and rectum; no polyps, diverticula, inflammation, or masses.

\section*{Follow-up Tests and Procedures}
Repeat screening in 10 years if normal, or sooner (3-5 years) depending on the size and pathology of polyps removed.

\section*{References}
\begin{enumerate}
  \item MedlinePlus. Colonoscopy. U.S. National Library of Medicine; 2025.
  \item Current Medical Diagnosis and Treatment. McGraw-Hill; 2025.
\end{enumerate}

\clearpage
